% !TeX root = ../初等数论笔记.tex
% 由主文档通过 \input 载入；请编译主文档。
\setcounter{chapter}{0} % 保留原书章号。
\chapter{常用符号与结论索引}
\label{app:reference}

\section{常用符号}
\begin{description}
 \item[整数与整除] $\mathbb Z_{>0}$ 为正整数集；$a\mid b$ 表示 $b=ak$ 对某个整数 $k$ 成立；$\gcd(a,b)$ 与 $\operatorname{lcm}(a,b)$ 分别为最大公因数与最小公倍数；$v_p(n)$ 为非零整数 $n$ 的素因子 $p$ 的指数。
 \item[同余与剩余] $a\equiv b\pmod m$ 表示 $m\mid a-b$；$[a]_m$ 为剩余类；$a^{-1}\pmod m$ 仅在 $\gcd(a,m)=1$ 时存在；$\operatorname{ord}_m(a)$ 是可逆类的乘法阶。
 \item[算术函数] $\tau,\sigma,\varphi,\mu$ 分别表示约数个数、约数和、欧拉函数与莫比乌斯函数；$f*g$ 为狄利克雷卷积；$\mathbf1$ 为恒一函数，$\varepsilon$ 为卷积单位元。$\sum_{d\mid n}$ 对正约数求和。
 \item[平方剩余] $\left(\frac ap\right)$ 在奇素数分母时为勒让德符号；推广到奇正分母时为雅可比符号，符号等于 $1$ 的可解性含义随分母类型改变。
 \item[连分数] $P_k/Q_k$ 为第 $k$ 个渐近分数，下标从 $0$ 开始；$\alpha_k$ 为完整商；平方根展开中的 $B_k,C_k$ 由式\eqref{eq:surd-recurrence}定义；$\ell$ 为最小周期长度。
 \item[分析记号] $\lfloor x\rfloor$ 与 $\lceil x\rceil$ 分别为下、上取整；$\pi(x)$ 为素数计数函数；$f\sim g$ 表示比值趋于 $1$，$f=O(g)$ 表示指定极限方向下的数量级上界。
\end{description}

\section{按问题查找结论}
\begin{description}
 \item[构造整数线性组合] 用贝祖等式（定理\ref{thm:bezout}），先排除两个输入同时为零的情形。
 \item[求一次整数方程] 用定理\ref{thm:linear-diophantine}；可解条件为最大公因数整除常数项，完整解见式\eqref{eq:linear-diophantine-solutions}。
 \item[比较素因子指数] 用算术基本定理（定理\ref{thm:unique-factorization}）及式\eqref{eq:factorial-valuation}。
 \item[求一次同余与同余组] 用定理\ref{thm:linear-congruence}、\ref{thm:crt}、\ref{thm:general-crt}；记录最终模数，不能把“一个解类”理解为“一个整数”。
 \item[求约数和或反演] 用式\eqref{eq:divisor-functions}、定理\ref{thm:totient}与\ref{thm:mobius-inversion}；积性只要求互素输入相乘时成立。
 \item[化简幂或求幂同余] 用费马、欧拉定理与阶的性质；使用指标前先核查模数是否有原根，见定理\ref{thm:primitive-root-classification}。
 \item[判断平方剩余] 用定理\ref{thm:euler-criterion}和\ref{thm:quadratic-reciprocity}；勒让德符号的分母必须为奇素数。
 \item[判断两平方和] 用定理\ref{thm:two-squares}；模 $4$ 余 $3$ 的素因子指数必须为偶数。
 \item[求 Pell 方程全部解] 先按定理\ref{thm:pell-fundamental}求基本解，再按定理\ref{thm:pell-all}生成；负方程使用命题\ref{prop:negative-pell}。
\end{description}

